\documentclass[11pt]{article}
\usepackage{amsmath,amssymb,amsthm}
\newtheorem{proposition}{Proposition}
\title{Weighted Flat Limits of Compact Reciprocal Sheets}
\author{}
\date{November 23, 2026}
\begin{document}
\maketitle

\paragraph{Status and scope.}
This note audits one attempted amplification mechanism for the reverse
compiler. It does not bound arbitrary admissible representations and
does not establish either requested asymptotic bound. The identities
and generic flat-limit statements below have written proofs; they are
not Lean formalized. The companion checker verifies 124 exact leading
Pl\"ucker coefficients for $H\leq5$.

Let $c_h,v_i\in\mathbb R^d$, let $v_{0i},\lambda_h\in\mathbb R$, and put
\[
 K_{ih}(\epsilon)=
 \frac1{1-\epsilon(v_{0i}+c_h\cdot v_i+\lambda_h)}.
\]
The tokens may be permuted among the columns. For any fixed finite
configuration, sufficiently small positive $\epsilon$ makes all these
denominators positive. Whole-cube implementation imposes the additional
affine slope budget, as in the earlier sheet constructions.

\begin{proposition}[Exact weighted expansion]
Write $U_i=1-\epsilon v_{0i}$ and $V_h=1-\epsilon\lambda_h$.
For sufficiently small $\epsilon$,
\[
 K_{ih}(\epsilon)=
 \sum_{\alpha\in\mathbb N^d,\ b\geq0}
 \binom{|\alpha|+b}{\alpha_1,\ldots,\alpha_d,b}
 \epsilon^{|\alpha|+2b}
 \frac{v_i^\alpha v_{0i}^b}{U_i^{|\alpha|+b+1}}
 \frac{c_h^\alpha\lambda_h^b}{V_h^{|\alpha|+b+1}}.
 \tag{1}
\]
The same identity holds as a formal power series.
\end{proposition}
\begin{proof}
Set $X=\epsilon(c_h\cdot v_i)/(U_iV_h)$ and
$Y=\epsilon^2v_{0i}\lambda_h/(U_iV_h)$. Then
\[
 U_iV_h(1-X-Y)
       =1-\epsilon(v_{0i}+c_h\cdot v_i+\lambda_h).
\]
Expand $(1-X-Y)^{-1}$ by the geometric series and expand powers of
$c_h\cdot v_i$ by the multinomial theorem. Multiplication by
$(U_iV_h)^{-1}$ gives (1). Each coefficient involves finitely many
indices, so the formal-series identity is equally valid.
\end{proof}

Thus parameter monomials $c^\alpha\lambda^b$ and row monomials
$v^\alpha v_0^b$ have weight $|\alpha|+2b$. Cauchy--Binet applied to
(1) gives an exact rule for the first possible coefficient of every
$H\times H$ minor: sum over all sets of $H$ distinct weighted indices
with minimum total weight, taking the product of the row evaluation
determinant and the parameter moment determinant, including the
multinomial coefficients. Higher-weight sets cannot contribute to
this coefficient. If it vanishes, the first actual nonzero coefficient
may occur later; the statements below explicitly exclude that case.

\paragraph{Second-order graph and its normalization.}
For $d=1$, $c=(0,1,2)$, and $H=3$, the leading column space is, whenever
the corresponding evaluations are independent,
\[
 \operatorname{span}\left\{1,\ v,\
 2v^2+(\lambda_0-2\lambda_1+\lambda_2)v_0\right\}.
 \tag{2}
\]
Indeed the only minimum-weight sets are
$(1,c,c^2)$ and $(1,c,\lambda)$, with total weight $3$.
Their parameter determinants are $2$ and
$\lambda_0-2\lambda_1+\lambda_2$, respectively.
More generally, if $a$ satisfies $\sum_h a_h=\sum_h a_hc_h=0$,
the second-order quotient is
\[
 \sum_h a_hc_h^2v^2+(a\cdot\lambda)v_0.
\]
The coefficient of $v_0$ is not $2(a\cdot\lambda)$: the raw expansion
also has an order-$\epsilon$ constant $a\cdot\lambda$, and cancelling
that constant by an $\epsilon$-dependent column combination subtracts
one copy of $(a\cdot\lambda)v_0$ at the next order.

For $H=4$, the first possible minor coefficient has order $5$ and
factors through the common space $\operatorname{span}\{1,v,v^2,v_0\}$.
For $H=5$, its order is $8$; the two boundary features are $v^3$ and
$2vv_0$, paired with moments $c^3$ and $c\lambda$. These formulas
are included in the exact checker.

\paragraph{Generic ambient dimension for one varying coordinate.}
Let $N_1(k)$ denote the number of pairs $(a,b)$ with $a+2b\leq k$.
Then
\[
 N_1(k)=\left\lfloor\frac{(k+2)^2}{4}\right\rfloor.
\]
Choose $k$ with $N_1(k-1)<H\leq N_1(k)$. Suppose a permutation has
at least one nonzero minor coefficient at the first possible total
weight. Every leading Pl\"ucker coordinate then factors through
the fixed row-evaluation space of weighted polynomials of degree at
most $k$. Consequently its flat column space is contained in this
common space, whose dimension is at most
\[
 N_1(k)=H+O(\sqrt H).
\]
This holds simultaneously for every permutation satisfying the
stated nonvanishing condition. A rank-deficient exceptional permutation
may have a later first nonzero coefficient and is not covered.

\paragraph{A bound independent of the coordinate dimension.}
The same first-layer argument gives a useful $2H$ bound even when
$d$ grows. Let $N_d(k)$ count weighted indices
$|\alpha|+2b\leq k$, and take
$N_d(k-1)<H\leq N_d(k)$. Assume again that the first possible
Pl\"ucker coefficient is nonzero. This requires full rank of all
lower-weight parameter moments. In particular, the pure moments
$c^\alpha$, $|\alpha|<k$, have rank
\[
 P_{k-1}=\binom{d+k-1}{d}<H.
\]
In a leading column-space vector, the pure degree-$k$ coefficient
is the image of a vector in the fixed kernel of those lower pure
moments. The image has dimension at most $H-P_{k-1}$ and does not
depend on the token permutation. Together with all lower pure row
monomials, this gives a fixed pure-polynomial space of dimension at
most $H$. All remaining row monomials contain $v_0$; dividing by
$v_0$ identifies them with weighted indices of weight at most $k-2$.
There are $N_d(k-2)<H$ of these. Therefore all such generic flat
spaces lie in one common space of dimension less than $2H$.

For clarity, the assertion about the boundary coefficient also follows
directly from the leading wedge in (1). Every minimum-weight index set
contains every lower-weight index. The leading column space contains
their row evaluations. Its remaining boundary vectors come from
parameter coefficient vectors annihilating all lower moments, hence
annihilating the stated pure moments. Projecting boundary vectors to
the pure degree-$k$ coordinates gives the fixed kernel image above.

\paragraph{What this does and does not rule out.}
These common ambient spaces cannot themselves realize all strict sign
patterns on more coordinates than their dimension. This is a statement
about the flat limits. A conclusion for small nonzero $\epsilon$ needs
a robustly missed orthant, for example an annihilating vector with
every coordinate nonzero. Deliberately degenerate query rows can fail
that condition: higher-order terms may determine signs at coordinates
that vanish in a leading vector. Likewise, exceptional permutations
can have deficient weighted moment matrices and later leading jets.
No recursive construction exploiting either possibility, and no
general obstruction to such a construction, is proved here.
\end{document}
